\section{Introduction}\label{sec:intro}
Location-based services (LBSs) connect private movement to remote information retrieval. A
device may request nearby facilities repeatedly, and a provider can combine the resulting
coordinates, timestamps and replies with road geometry and prior knowledge. Correlation across
requests matters: a plausible replacement at one instant may become identifiable when the
complete sequence is examined. The uniqueness of mobility traces motivates this concern,
although observations on real populations do not automatically transfer to synthetic
trajectories or to physical identities~\cite{mobilityunique}.

Geo-indistinguishability (Geo-I) expresses location privacy through a likelihood-ratio bound
scaled by input distance~\cite{geoi}. Extending that guarantee to an implemented continuous
service requires accounting for more than repeated noise generation. Decisions to reuse a
reference, the construction of query candidates, the sampling of public actions, retries and
session boundaries can each depend on private information. Predictive private mechanisms
already show the importance of managing new information over time~\cite{predictive}.
Our focus is the explicit connection between these mechanisms and the transcript actually
observed by an LBS provider.

Utility also requires a task definition. A coordinate close to a user may retrieve poor
answers across a one-way road or a travel-time boundary. Conversely, several remote query
locations can recover useful metadata that the device ranks accurately using its current
GPS. Prior work combines Geo-I, dummy queries and local processing in application-specific
settings~\cite{ev}. Thus neither multiple queries nor local selection alone is our novelty
claim. The question is how to specify a stateful publication policy, bound its complete
transcript, and evaluate its answers without confusing a planning proxy with service utility.

Comparability is another obstacle. Recent methods release a replacement coordinate, a set
containing the true location and dummies, or extra fake packets; offline trajectory generation
has a different output contract~\cite{transprotect,semantic,fake,sok}. A real-member anonymity
metric does not apply unchanged to a query set that need not contain the true location.
We therefore retain task-aligned localization and service metrics alongside original metrics
where their definitions and required inputs remain applicable. We do not propose new formulas
for mean absolute error (MAE), Recall or bootstrap intervals.

We examine \BQ{}, a family of policies with Geo-I cap $C_s$ and query allowance $\Gamma_Q$;
the evaluated v1 configuration remains fixed throughout this article. A road exponential
mechanism and a private reuse test maintain an internal reference $Z$. An approximate belief
computed from protected history scores a common library of road-feasible query segments.
A bounded centered kernel selects one segment; its frames are then released at public times.
The device retrieves a purpose-independent cover, caches valid metadata, and returns a full
filtered and sorted local answer.

The information-access model is consequential. The present prototype loads a complete
public reference that matches the static provider used in the experiment. Its utility
evaluation restricts answers to received candidates, but does not establish that an
operational device lacks the information to answer offline. Likewise, a nonadaptive public
base has slightly higher Recall in this cohort. We preserve these adverse findings.
They limit claims about application necessity and the benefit of private adaptation, while
leaving the conditional transcript analysis meaningful.

The article makes three scoped contributions:
\begin{enumerate}[leftmargin=*,itemsep=3pt]
\item A road-aware, stateful publication policy with separate accounting for private GPS
operations and bounded joint query-segment selection. The analysis covers the ordered
transcript under stated support, clock, provider and state assumptions.
\item A connected privacy--utility analysis: an end-to-end likelihood-ratio bound, a
prior-concentration corollary for arbitrary-estimator localization, conditional composition
of multiple channels, and utility statements that separate exact properties from
uncertified calibration assumptions.
\item A reproducible controlled comparison with three audited paper adaptations and two
controls, accompanied by supplementary inference tasks, family-level uncertainty, cost
accounting and an audit of catalogue saturation and the untested cap-exhausted regime.
\end{enumerate}
We do not claim to be the first road-aware Geo-I mechanism, to outperform every original
paper implementation, or to certify a complete deployed system.
Section~\ref{sec:related} positions the work; Section~\ref{sec:method} specifies and analyzes
the method; Section~\ref{sec:evaluation} reports the comparison and its limits; and
Section~\ref{sec:conclusion} concludes.
